\section{Discussion: why detecting is not causing here}
\label{sec:discussion}

\paragraph{A marker of a state is not the cause of an action.} Looping is a property of the trajectory: the agent repeats because the screen did not change or the target was misread. A direction that encodes ``this episode is stuck'' can be read reliably while the repeated action is computed from features that do not depend on it. The global-erasure control shows this directly. Probes in computer-use agents can reach very high AUC without detecting the intended property \citep{li2026auc}, and format-confounded probes can lack any functional role \citep{sahoo2026format}. Our 0.8B result is a case where the probe does track the intended property and still offers no lever.

\paragraph{False success looks like a belief that the task is done, acted on without checking.} The 4B evidence points one way. The model says ``done'' on most false completions as on true ones, and its reasoning writes the conclusion first in both cases, a property of its output format (\S\ref{sec:think}). The tool choice becomes visible late, at the tool name, after a thought that already declares success, through a circuit that reads the tool definition and treats true and false completions the same (\S\ref{sec:map}). Pushing toward low value makes the model write that it should check, without making it check. This matches reports that stated confidence tracks commitment more than correctness \citep{kumaran2026commit} and that stopping decisions are spread over late components \citep{dutta2026termination}. The perceived-completion direction (\S\ref{sec:perceive}) adds that, on \pcMatchN{} matched pairs, its projection reads false completions as complete. The public report of an agent that claims success after acting with no screenshot at all \citep{holodesktop2026issue15} is the extreme case of the same pattern: the claim of completion is not conditioned on evidence that the task was done. Whether that state is a misperception of the screen or a commitment to finish, removing it is likely to suppress true completions too, and the \fsm{} criterion of our gate exists to catch exactly that. The clearest single case argues against a pure perception error. On a Chrome rental task, the model reads the current date ``Nov 8'' on screen, then states that ``Nov 10--11'' is next month and answers ``done''; a random direction of the same norm produces the same contradiction. % source: data/pod_astra/verify/semantic_reviews.json (verify_a1 sample 3: "Explicitly reads current date Nov8 then calls Nov10–11 next month"); data/pod_astra/verify/REPORT.md
The model looked; it compared wrongly. A second sample on the same task, under the ``Wait'' edit, goes further: it derives that next month is December, says the dates must be changed, then decides that the main requirement (sorting by seats) is met and answers ``done''. % source: data/pod_astra/mit1/REPORT.md "New mechanistic observation under wait2"
One output is not a population claim, but it suggests that a false ``done'' can be the partial satisfaction of a conjunctive task, with the unmet requirement noticed and then dropped. A repair may therefore have to enforce that every requirement holds, rather than act on looking or on the decision to stop.

\paragraph{Does the agent know?} An external reader asked whether an agent that knows a requirement is still unmet answers ``done'' anyway. Under a protocol registered before annotation, a blind annotator read the instruction and the note and thought of the last three steps of every episode in which Holo-3.1-4B ends with \texttt{answer} on the \textsc{test} split, in both servings of the original (agreement \cdAudit{} on an independent re-annotation). When it ends with \texttt{answer}, the agent rarely writes in those steps that a requirement remains unmet: in \cdFs{} false successes (\cdFsPct\%) against \cdTs{} true successes (\cdTsPct\%), a difference of \cdDiff{} points (95\% CI \cdDiffCi; one-sided Fisher $p=\cdP$). Most of these mentions state that the task cannot be done, and \cdRegex{} of the nine on the false-success side are correct infeasibility answers that our harness scored as ``done''. Dropping a known constraint therefore explains few false successes (two or three of the 56), and a free rule based on this signal would catch only \cdFsPct\% of them, against \gbClaudeDet\% for the external verifier (on a different set of episodes, so the comparison is indicative). % source: data/local/constraint_drop/RESULT.json, REPORT.md ("2 or 3 of the 56"), PROTOCOL.md (commits b88886d, 42a99b0) (figs/guard.py)

\paragraph{Implications for agent developers.} We found no direction whose removal lowers false ``done'' selectively, in screens able to detect only large effects; within our probes, nothing separated false from true completions beyond the episode's length except the hesitation between samples (Table~\ref{tab:fs-auc}, \S\ref{sec:levers}), and no edit that moved the decision showed an established selectivity on feasible tasks: the merge showed none, and the closed-form edit's interval includes zero. Our results suggest that a fix must bring in information the model does not use when it stops: an external verifier that checks each requirement of the task against the screen before \texttt{answer} is accepted (\S\ref{sec:guard}: offline, it separates false from true successes where the agent, judging itself with a free-form prompt, does not; online, it lowers the false ``done'' rate without a measurable score change), or training that targets false successes, for instance with rewards that penalise episodes ending in a false ``done''. Weight surgery of the kind tested here has not shown that it can supply it.

\paragraph{What would change our mind.} The official analysis has already ruled out the first condition we set, a direction on at least 40 \fsp{} steps that beats position (Table~\ref{tab:fs-auc}). What remains is an intervention that moves \fsp{} answers while sparing \fsm{} on unseen tasks, perhaps at the larger doses the ceiling effect calls for, and a paired online gain on the confirmation reserve larger than the \onMde{} points our reserve can detect. Measuring how many false-success states can still be repaired when \texttt{answer} is forbidden would bound what any such edit can gain; we have not measured it. % source: data/local/reports/REDTEAM_0305.md point 2
